% ======================================================================
% 10_design_control_vision.tex: CHAPTER 3, design and implementation
% WHAT TO WRITE
%   Requirements, overall architecture, then each part in detail (hardware, software, algorithms, parameters).
%   This is usually the longest chapter: use many \subsection and \paragraph headings to keep it readable.
%   Show block diagrams as figures and parameter lists as tables.
%   "Side-by-Side Figures" shows two related pictures next to each other with a caption each.
%
% HOW TO USE THIS FILE
%   - Replace every ... with your text. A new paragraph starts after one blank line.
%   - The title in \chapter{...} is printed as "CHAPTER n: TITLE" in capitals. Rename it freely.
%   - \section{...} and \subsection{...} are numbered automatically (1.1, 1.1.1). Rename, copy or delete them.
%   - \paragraph{Small heading:} gives a bold line of its own (the colon is required).
%   - The last section ("Template Examples") shows a figure, an equation and a table. Copy a block into your
%     text, edit it, then delete the "Template Examples" section when you no longer need it.
%   - Syntax for everything (lists, figures, tables, equations, citations): docs/LATEX_CHEATSHEET.md
%   - Labels (\label{...}) must be unique across the whole report.
% ======================================================================

\chapter{SYSTEM DESIGN AND IMPLEMENTATION}
\label{chap:design}

\section{System Requirements}
...

\subsection{Functional requirements}
...

\subsection{Technical specifications}
...

\section{Architecture}
...

\subsection{Overall architecture}
...

\subsection{Block diagram}
...

\section{Detailed Design}
\subsection{Component A}
...

\paragraph{Sub-topic one:}
...

\paragraph{Sub-topic two:}
...

\subsection{Component B}
...

\subsection{Side-by-Side Figures}
Two related figures placed side by side (Figure~\ref{fig:side-a} and Figure~\ref{fig:side-b}):
\begin{figure}[htbp]
    \centering
    \begin{subfigure}[b]{0.48\linewidth}
        \centering
        \includegraphics[width=\linewidth]{figures/example_a.png}
        \caption{...}
        \label{fig:side-a}
    \end{subfigure}\hfill
    \begin{subfigure}[b]{0.48\linewidth}
        \centering
        \includegraphics[width=\linewidth]{figures/example_b.png}
        \caption{...}
        \label{fig:side-b}
    \end{subfigure}
    \caption{...}
    \label{fig:side-by-side}
\end{figure}

\section{Implementation}
...

\subsection{Hardware or software setup}
...

\subsection{Procedure}
...

% ============================================================
% Template examples: copy these blocks into the relevant subsection.
% Replace placeholder text, captions, labels, and sample data as needed.
% ============================================================
\section{Template Examples}
\label{sec:template-ch3-design-control-vision}

\subsection{Figure and Cross-Reference Example}
...

Figure~\ref{fig:template-ch3-design-control-vision} illustrates ...
\begin{figure}[htbp]
    \centering
    % Replace the framed placeholder with your image, for example:
    % \includegraphics[width=0.75\linewidth]{figures/example.png}
    \fbox{%
        \begin{minipage}[c][4cm][c]{0.75\linewidth}
            \centering
            ...
        \end{minipage}%
    }
    \caption{...}
    \label{fig:template-ch3-design-control-vision}
\end{figure}

\subsection{Equation Example}
...

An example relationship is given in Equation~\eqref{eq:template-ch3-design-control-vision}:
\begin{equation}
    y(t) = a\,x(t) + b,
    \label{eq:template-ch3-design-control-vision}
\end{equation}
where $x(t)$ is the input, $y(t)$ is the output, and $a$ and $b$ are
example coefficients. Replace this relationship and its definitions with ...

\subsection{Table Example}
...

Table~\ref{tab:template-ch3-design-control-vision} presents ...
\begin{table}[htbp]
    \centering
    \caption{...}
    \label{tab:template-ch3-design-control-vision}
    \begin{tabularx}{0.9\linewidth}{lXX}
        \toprule
        \textbf{Item} & \textbf{Description} & \textbf{Value / Unit} \\
        \midrule
        ... & ... & ... \\
        ... & ... & ... \\
        \bottomrule
    \end{tabularx}
\end{table}
